% ============================================================
% Updated Evidence: 483 Text Signals and Patient-Harm Risk
% Rerun of the INFORMS 2026 VAI-only result with the current,
% human-validated extraction pipeline (Claude Sonnet 5, v2 prompt)
% Prepared 2026-09-16
% ============================================================
\documentclass[aspectratio=169,10pt]{beamer}

% -- Theme -------------------------------------------------------------------
\usetheme{Madrid}
\usecolortheme{default}
\setbeamertemplate{navigation symbols}{}
\setbeamertemplate{footline}[frame number]
\setbeamersize{text margin left=6mm, text margin right=6mm}

% -- Packages ------------------------------------------------------------------
\usepackage{amsmath, amssymb}
\usepackage{graphicx}
\usepackage{booktabs}
\usepackage{array}
\usepackage{tabularx}
\usepackage{xcolor}
\usepackage{colortbl}
\usepackage{hyperref}
\pdfstringdefDisableCommands{\let\\\space\let$\relax}

% -- Colors -- matches Paper/Text Analysis/20260714_informs_slides.tex --------
\definecolor{ncsuRed}{RGB}{204,0,0}
\definecolor{ncsuGray}{RGB}{100,100,100}
\definecolor{techBlue}{RGB}{37,99,235}
\definecolor{govGreen}{RGB}{5,150,105}
\definecolor{warnOrange}{RGB}{217,119,6}
\definecolor{lightGray}{RGB}{243,244,246}

\setbeamercolor{structure}{fg=ncsuRed}
\setbeamercolor{block title}{bg=ncsuRed!90, fg=white}
\setbeamercolor{block body}{bg=ncsuRed!8}
\setbeamercolor{alerted text}{fg=ncsuRed}

\newcommand{\highlight}[1]{{\color{ncsuRed}\textbf{#1}}}
\newcommand{\tech}[1]{{\color{techBlue}\textbf{#1}}}
\newcommand{\gov}[1]{{\color{govGreen}\textbf{#1}}}

% -- Title ---------------------------------------------------------------------
\title[Updated Text-Signal Evidence]{%
  Updated Evidence: 483 Text Signals\texorpdfstring{\\}{ }
  and Patient-Harm Risk\texorpdfstring{\\}{ }
  \small A rerun of the INFORMS VAI-only result on the current, validated pipeline}

\author[Sahebi-Fakhrabad et al.]{Amirreza Sahebi-Fakhrabad}
\date{September 16, 2026 -- internal update, not for external distribution}

\begin{document}

% ─────────────────────────────────────────────────────────────────────────────
\begin{frame}
\titlepage
\end{frame}

% ─────────────────────────────────────────────────────────────────────────────
\begin{frame}{Why This Update}
\begin{itemize}
  \setlength{\itemsep}{10pt}
  \item INFORMS Annual Meeting (July 2026) presented a flagship result: 483 text signals
        predict future patient-harm risk in facilities where FDA's own classification (VAI)
        carries the least information.
  \item That analysis used an earlier extraction (before Claude Sonnet 5 was pinned as the
        validated model, and before the round-1 human-label validation and prompt fixes
        documented in the LLM Extraction Validation Report).
  \item \highlight{This deck reruns the exact same test on the current, validated data} --
        same panel construction, same cohort definition, same models -- to see whether the
        claim holds up.
\end{itemize}
\end{frame}

% ─────────────────────────────────────────────────────────────────────────────
\begin{frame}{Study Design (unchanged from INFORMS)}
\begin{itemize}
  \setlength{\itemsep}{8pt}
  \item \textbf{Unit:} one inspection event
  \item \textbf{Features:} 17 LLM text signals from the 483 (now from the validated v2 prompt,
        Claude Sonnet 5)
  \item \textbf{Outcome:} above-median adverse events in the 4 quarters after inspection
  \item \textbf{CV:} FEI-grouped 5-fold logistic regression and random forest
  \item \textbf{Cohort:} inspections with an ANDA-matched FAERS record in the outcome window --
        \highlight{176 inspections, 78 FEIs} (98 FEIs, 33 OAI, before this filter)
  \item \textbf{VAI-only subgroup:} 55 facilities that never received an OAI in this panel --
        \highlight{110 inspections}
\end{itemize}
\vspace{6pt}
{\small Cohort sizes above match the original INFORMS slide exactly -- confirms this rerun is a
faithful replication of the original setup, not a different test.}
\end{frame}

% ─────────────────────────────────────────────────────────────────────────────
\begin{frame}{Sanity Check: The Structured-Data-Only Config}
\begin{columns}[T]
\column{0.55\textwidth}
\begin{block}{Config C: OAI flag only (no text)}
\begin{itemize}
  \item Does not use any 483 text.
  \item Should be unaffected by anything that changed in the text pipeline.
\end{itemize}
\end{block}
\column{0.4\textwidth}
\begin{center}
{\small INFORMS (July):} \highlight{0.545}\\[4pt]
{\small Rerun (Sept):} \highlight{0.537}
\end{center}
\end{columns}
\vspace{10pt}
Matches closely. This is strong evidence the rerun's panel construction, FAERS matching, and
cross-validation are a faithful reproduction of the INFORMS methodology -- so where the
\emph{text}-dependent configs differ below, that is a real finding, not a plumbing difference.
The small residual gap (0.545 vs.\ 0.537) is from an unrelated data refresh (the panel was
rebuilt against updated FDA Inspection data), not anything text-related.
\end{frame}

% ─────────────────────────────────────────────────────────────────────────────
\begin{frame}{Updated Results}
{\small
\begin{tabular}{clcc}
\toprule
 & Feature set & INFORMS (July) & Current (Sept 16, corrected feature) \\
\midrule
A & Text only, full sample (LR) & 0.585 & \highlight{0.483} \\
C & OAI flag only & 0.545 & 0.537 \\
\midrule
\pause
D & Text only, VAI-only (LR) & 0.656, $p{=}0.046$ & \highlight{0.591, $p{=}0.189$} \\
D & Text only, VAI-only (RF) & not reported & \highlight{0.610, $p{=}0.080$} \\
\bottomrule
\multicolumn{4}{l}{\scriptsize $p$ from one-tailed $t$-test vs.\ 0.5 across 5 CV folds, same test as the original slide}
\end{tabular}
}
\vspace{10pt}
\pause
\begin{alertblock}{Honest read}
\small The flagship VAI-only claim does \textbf{not} replicate at the original strength under
logistic regression, the exact model INFORMS reported (p goes from 0.046 to 0.189). The random
forest version, which had looked significant in an earlier rerun of this analysis (p=0.032),
moved further from significance once a severity-feature bug fix landed the same day (p=0.080).
Neither model clears p$<$0.05 now. Both still point the same direction, so there is a real signal,
but it is suggestive, not statistically established.
\end{alertblock}
\end{frame}

% ─────────────────────────────────────────────────────────────────────────────
\begin{frame}{Why the Numbers Moved}
\begin{itemize}
  \setlength{\itemsep}{10pt}
  \item The text features themselves changed twice: first, the current extraction is
        human-validated (independent blind review, 68--96\% field accuracy) with two real prompt
        bugs and several rule gaps fixed after the INFORMS numbers were produced. Second, a
        severity-feature bug fix landed the same day this rerun was first drafted:
        severity\_tier's real confusion is at the Major/Moderate boundary, not Critical/Major, so
        severity\_majmod\_share replaced severity\_critmajor\_share as the severity feature
        everywhere, moving the VAI-only numbers again (RF: p=0.032 $\to$ p=0.080).
  \item This is not a case of ``the old number was wrong, ignore it'' -- it is the expected
        result of tightening the pipeline twice in the same direction. A weaker, no-longer-
        significant effect on more rigorously validated data is a more honest basis to build on
        than a stronger effect on data known to have had extraction bugs at the time.
  \item Sample size remains the binding constraint either way: 55 VAI-only facilities, 110
        inspections. INFORMS's own slide called this subgroup result ``hypothesis-generating,''
        and that framing applies at least as much now.
\end{itemize}
\end{frame}

% ─────────────────────────────────────────────────────────────────────────────
\begin{frame}{Where This Points Next}
\begin{itemize}
  \setlength{\itemsep}{9pt}
  \item \textbf{The shortage test has now been run (m19), with a significance test.} Direct test
        (text-only, VAI-only, shortage as outcome): mixed -- logistic not significant (AUC 0.514,
        $p{=}0.243$), Random Forest significant (AUC 0.590, $p{=}0.013$) despite only 28 events.
        Worth taking seriously, not dismissing as too sparse.
  \item \textbf{An unrelated, more promising finding from the same model.} A plain inspection +
        structural baseline (no text, full 127-facility universe) shows real signal for shortage
        exposure and clears significance comfortably: AUC 0.540--0.610, $p{=}0.003$ and
        $p{=}0.011$ -- evidence for the broader claim (quality/regulatory signals predict
        shortage), though it doesn't confirm the text-specific VAI hypothesis above. See
        RESULTS.docx.
  \item \textbf{More data, not a different model.} The VAI-only cohort is still small (55
        facilities) because 483-text coverage is still short for many facilities. Sample size,
        not signal quality, is still what keeps the logistic-regression versions from clearing
        significance. Revisit as coverage grows.
  \item \textbf{Done: OAI/VAI source aligned with m19.} This analysis now matches m19's source
        order (FDA Inspection Details primary, Redica fallback). Rerunning reproduced the
        identical cohort and identical AUCs -- a useful robustness check, not a source of the
        movement seen on this deck.
\end{itemize}
\end{frame}

\end{document}
